\documentclass{article}
\usepackage[margin=0.8in]{geometry}
\usepackage{hyperref}

\setlength{\parskip}{5pt}

\title{Theoretical and Mathematical Foundations\\of Spacecraft ADCS}
\author{Rishav}
\date{January 25, 2023}

\begin{document}
\maketitle
Last modified: \today

\newpage
\section*{Preface}
On October 26, 2019, I was introduced to Attitude Determination and Control System (ADCS) as a potential topic for my final year engineering project by Rakesh Chandra Prajapati. I spent the next three years, for and after my thesis, studying and implementing the ideas and algorithms in ADCS. This project is my attempt to document the experience with the hope that it will act as a beginner's guide to ADCS for those who want to pursue the field seriously.

No prior prerequisites are assumed to read and understand the contents. I hope you enjoy it.

\newpage
\tableofcontents

\newpage
\section{Introduction}
\subsection{Etymology of \textit{attitude}}
\subsection{ADCS of spacecrafts}

\section{Measurements}

\section{Frame of references}
\subsection{Body frame}
\subsection{Inertial and non-inertial frames}
\subsection{Transformation between different frames}

\section{Reference frames in ADCS}
\subsection{Earth Centered Inertial Frame (ECI)}
\subsection{Earth Centered Earth Fixed Frame (ECEF)}
\subsection{North-East-Down frame (NED)}
\subsection{Perifocal frame}

\section{Attitude parameterization}
\subsection{Euler angles}
\subsection{Rotation matrices}
\subsection{Quaternions}

\section{Attitude determination}
\subsection{Wahba's Problem}
\subsection{Rotation matrix}
\subsection{Quaternions}
\subsection{Sensor and model noises}

\section{Attitude estimation}
\subsection{Linear Kalman filter}
\subsection{Extended Kalman filter}
\subsection{Multiplicative Extended Kalman filter}

\newpage
\appendix
\section{Differential equations}

\end{document}
